% =====================================================================
% Module 06 — Dual Compatibility
% Academic-Galactic v1.2
% =====================================================================

\section{Dual Compatibility}
\label{sec:DualCompatibility}

\subsection{Overview}
Dual compatibility defines how dual operators act on the Soft 80-dimensional
manifold $\mathcal{M}^{80}$ and ensures consistency between primal and dual
structures across filtration layers, fiber bundles, spectral bundles, and
collapse surfaces.

The dual operator is denoted:
\[
\mathcal{C} : \mathcal{F}_k \to \mathcal{F}_k,
\]
and interacts with:
\begin{itemize}
    \item spectral operators $\Phi_5$ (Module 03),
    \item curvature envelopes $\Phi_7$ (Module 02),
    \item holonomy transport $\Phi_8$ (Module 02),
    \item topology-modulation drift $\Phi_8^t$ (Module 02),
    \item completion operator $\Phi_9$ (Module 08),
    \item reprojection operator $\mathcal{R}$ (Module 07),
    \item domain and altitude operators (Module 01),
    \item composite operator algebra (Module 05).
\end{itemize}

Dual-operator foundations follow operator-algebraic state-space geometry
\citeClassA{AlfsenSchultz2003} and categorical-homotopical operator algebra
\citeClassD{EmmersonBuchanan2025}.

\subsection{Dual Operator Definition}
\begin{definition}[Dual Operator]
A dual operator at altitude $k$ is a cotangent-compatible map
\[
\mathcal{C} : \mathcal{F}_k \to \mathcal{F}_k
\]
satisfying:
\[
\mathcal{C}^2 = \mathrm{id},
\]
and preserving cotangent structure.
\end{definition}

Dual operators act on both primal and dual bundles.

\subsection{Cotangent Compatibility}
Cotangent compatibility requires:
\[
\mathcal{C}(T_p\mathcal{F}_k) = T_p^\ast\mathcal{F}_k,
\]
ensuring that dual operators map tangent structure to cotangent structure.

This compatibility is altitude-preserving:
\[
\mathcal{C}(\mathcal{F}_k) \subseteq \mathcal{F}_k.
\]

\subsection{Spectral–Dual Alignment}
Spectral operators commute with dual operators:
\[
[\Phi_5, \mathcal{C}] = 0.
\]

This ensures:
\begin{itemize}
    \item cotangent spectral consistency,
    \item spectral invariant stability,
    \item compatibility with spectral curvature (Module 03).
\end{itemize}

\subsection{Curvature–Dual Interaction}
Curvature envelopes satisfy:
\[
\mathcal{C}(\Phi_7(p)) = -\Phi_7(\mathcal{C}(p)),
\]
indicating that curvature is dual-antisymmetric.

This dual antisymmetry is consistent with:
\begin{itemize}
    \item holonomy-curvature interaction \citeClassA{BlackholeHolonomy2026},
    \item informational holonomy curvature \citeClassA{GutierrezUle2026},
    \item discrete-to-continuous curvature convergence \citeClassA{LiPengDong2026}.
\end{itemize}

\subsection{Holonomy–Dual Compatibility}
Holonomy transport satisfies:
\[
\mathcal{C}(\Phi_8(\gamma)) = \Phi_8(\gamma^{-1}),
\]
where $\gamma^{-1}$ is the reversed curve.

This ensures:
\begin{itemize}
    \item dual holonomy consistency,
    \item compatibility with collapse-surface drift (Module 04),
    \item stability under topology modulation.
\end{itemize}

\subsection{Drift–Dual Interaction}
Topology-modulation drift satisfies:
\[
[\Phi_8^t, \mathcal{C}] \neq 0,
\]
indicating that drift induces dual-asymmetric topology change.

This dual asymmetry is essential for:
\begin{itemize}
    \item collapse-surface evolution,
    \item stratified topology modulation,
    \item reprojection consistency (Module 07).
\end{itemize}

\subsection{Reprojection–Dual Consistency}
Reprojection operator $\mathcal{R}$ satisfies:
\[
\mathcal{C}(\mathcal{R}(p)) = \mathcal{R}(\mathcal{C}(p)),
\]
ensuring dual-stable geometric consistency restoration.

\subsection{Completion–Dual Stability}
Completion operator $\Phi_9$ satisfies:
\[
[\Phi_9, \mathcal{C}] = 0,
\]
ensuring that invariant extraction is dual-stable.

\subsection{Summary}
Dual compatibility defines cotangent structure, spectral–dual alignment,
curvature–dual antisymmetry, holonomy–dual consistency, drift–dual interaction,
and reprojection–dual stability. These structures ensure that dual operators
interact coherently with the composite operator algebra and the geometric and
spectral structures of the Soft 80D manifold.

% =====================================================================
% End Module 06
% =====================================================================
